\documentclass[a4paper,11pt]{article}
\usepackage[T1]{fontenc}
\usepackage[utf8]{inputenc}
\usepackage{lmodern}
\usepackage{xcolor}
\usepackage{textcomp}
%\usepackage{newspaper}
%\date{\today}
%\currentvolume{1430}
%\currentissue{6}
%\usepackage{times}
\usepackage{graphicx}
\usepackage{multicol}
\usepackage{etoolbox}
\usepackage{picinpar}
%uasage of picinpar:
%\begin{window}[1,l,\includegraphics{},caption]xxxxx\end{window}
\usepackage{times}
\usepackage[T1]{fontenc}
\usepackage[utf8]{inputenc}
%\usepackage[nohead, nomarginpar, margin=1in, foot=.25in]{geometry}
\usepackage{lmodern}
%\usepackage{fontspec}
\usepackage{setspace}
%\usepackage{fancyhdr}
\usepackage{everyshi}
\usepackage{pdfpages}
\usepackage{indentfirst}
\usepackage[hyphens]{url}
\usepackage{hyperref}
\usepackage{pdfpages}
\usepackage{pgfpages}
\usepackage{amssymb}
\usepackage{amsmath}
\usepackage[notextcomp]{stix}
 \usepackage{paracol}
 \usepackage{stmaryrd}
 \usepackage{xurl}
 \usepackage{wasysym}
\usepackage[left=2cm, right=2cm, top=2cm, bottom=2cm]{geometry}
%\pagestyle{fancy}
%\lhead{Curriculum vitae, Anamitro Biswas}
%\chead{}
%\rhead{}
%\lfoot{}
%\cfoot{\thepage}
%\rfoot{}
%\makeatletter
%\renewcommand{\@chapapp}{Document}
%\makeatother
%\renewcommand\refname{\enhead\color{red}{References:}}
%\renewcommand{\headrulewidth}{0.4pt}
%\renewcommand{\footrulewidth}{0.4pt}
%\newfontfamily\bn{TiroBangla-Regular.ttf}[Script=Bengali]
%\newfontfamily\enhead{RubikDoodleShadow-Regular.ttf}[Script=Latin]
%\setmainfont{cmunrm.ttf}[Script=Latin]
%\newfontfamily\bn{TiroBangla-Regular.ttf}[Script=Bengali]
%\newfontfamily\latinfont{OfenbacherSchwabCAT.ttf}[Script=Latin]


%BRACKETS:
\newcommand{\lb}{\left(}
\newcommand{\rb}{\right)}
\newcommand{\lc}{\left\lbrace}
\newcommand{\rc}{\right\rbrace}
\newcommand{\ls}{\left[}
\newcommand{\rs}{\right]}

%GREEK UPPERCASE:
\newcommand{\Lm}{\Lambda}
\newcommand{\G}{\Gamma}
\newcommand{\D}{\Delta}
\newcommand{\Th}{\Theta}
\newcommand{\F}{\Phi}

%GREEK LOWERCASE:
\newcommand{\m}{\mu}
\newcommand{\lm}{\lambda}
\newcommand{\vs}{\varsigma}
\newcommand{\g}{\gamma}
\renewcommand{\t}{\tau}
\renewcommand{\a}{\alpha}
\renewcommand{\b}{\beta}
\newcommand{\e}{\varepsilon}
\newcommand{\dt}{\delta}
\newcommand{\f}{\varphi}
\newcommand{\z}{\zeta}
\newcommand{\s}{\sigma}
\newcommand{\p}{\pi}

%SET THEORY:
\newcommand{\cl}{\overline}
\newcommand{\intr}{^\circ}
\newcommand{\ir}{^\text{int}}
\newcommand{\comp}{^C}

%MATHCAL:
\newcommand{\A}{\mathcal{A}}
\newcommand{\B}{\mathcal{B}}
\newcommand{\C}{\mathcal{C}}
\newcommand{\cD}{\mathcal{D}}
\newcommand{\E}{\mathcal{E}}
\newcommand{\cF}{\mathcal{F}}
\newcommand{\cG}{\mathcal{G}}
\newcommand{\Hy}{\mathcal{H}}
\newcommand{\I}{\mathcal{I}}
\newcommand{\J}{\mathcal{J}}
\newcommand{\List}{\mathcal{L}}
\newcommand{\M}{\mathcal{M}}
\newcommand{\N}{\mathcal{N}}
\newcommand{\R}{\mathcal{R}}
\newcommand{\Pa}{\mathcal{P}}
\newcommand{\La}{\mathcal{L}}
\newcommand{\Ta}{\mathcal{T}}
\newcommand{\mcS}{\mathcal{S}}
\newcommand{\V}{\mathcal{V}}
\newcommand{\X}{\mathcal{X}}
\newcommand{\Y}{\mathcal{Y}}

%MATHFRAK:
\newcommand{\Mf}{\mathfrak{M}}
\newcommand{\Nf}{\mathfrak{N}}
\newcommand{\gf}{\mathfrak{g}}
\newcommand{\mfi}{\mathfrak{i}}
\newcommand{\mfj}{\mathfrak{j}}
\newcommand{\mfa}{\mathfrak{A}}

%FIELDS:
\newcommand{\reals}{\mathbb{R}}
\newcommand{\bC}{\mathbb{C}}
\newcommand{\bG}{\mathbb{G}}
\newcommand{\bP}{\mathbb{P}}
\newcommand{\bE}{\mathbb{E}}
\newcommand{\bH}{\mathbb{H}}
\newcommand{\bF}{\mathbb{F}}
\newcommand{\bN}{\mathbb{N}}
\newcommand{\naturals}{\mathbb{N}}
%MATH BOLD:
\newcommand{\bfa}{\mathbf{a}}
\newcommand{\bfb}{\mathbf{b}}
\newcommand{\bfg}{\mathbf{g}}
\newcommand{\bft}{\mathbf{t}}
\newcommand{\bfx}{\mathbf{x}}
\newcommand{\bfy}{\mathbf{y}}
\newcommand{\bfzero}{\mathbf{0}}

%GROUP THEORY:
\newcommand{\inv}{^{-1}}

%FUZZY TOPOLOGY:
\newcommand{\nphi}{\bigodot}

%ALGEBRAIC TOPOLOGY:
\newcommand{\rel}{\text{ rel}}

%HOMTOPY GROUPS:
\newcommand{\htopo}{\mathcal{H}\mathfrak{o}\mathcal{T}\mathfrak{o}\mathfrak{p}_*}
\newcommand{\topo}{\mathcal{T}\mathfrak{o}\mathfrak{p}_*}
\newcommand{\htop}{\mathcal{H}\mathfrak{o}\mathcal{T}\mathfrak{o}\mathfrak{p}}
\newcommand{\cattop}{\mathcal{T}\mathfrak{o}\mathfrak{p}}

%MAPS:
\newcommand{\tends}{\rightarrow}
\newcommand{\id}{\text{id}}


%DASH:
\newcommand{\dash}{^\prime}
\newcommand{\ddash}{^{\prime\prime}}
\newcommand{\tdash}{^{\prime\prime\prime}}


\newcommand{\Ob}{\text{Ob}}
\newcommand{\Eq}{\text{Eq}}




%HEADINGS:
\newtheorem{thm}{Theorem}
\newtheorem{prop}[thm]{Proposition}
\newtheorem{cor}[thm]{Corollary}
\newtheorem{defn}[thm]{Definition}
\newtheorem{obs}[thm]{Observation}
%\newtheorem{prop}[thm]{Proposition}
\newtheorem{lem}[thm]{Lemma}
\newtheorem{conj}[thm]{Conjecture}
\newtheorem{claim}[thm]{Claim}
\newtheorem{note}[thm]{Note}
\newtheorem{notation}[thm]{Notation}
\newtheorem{rem}[thm]{Remark}
\newtheorem{prob}[thm]{Problem}

%\setmainfont{Cormorant Garamond SemiBold}
\urlstyle{same}
\usepackage{tikzpagenodes}
\usepackage{tikz}
\usepackage{eso-pic}
\begin{document}
\title{Basic Homotopy Theory, and CW Complexes: How To Build Them}
\author{Anamitro Biswas\\Email: anamitroappu@gmail.com}
\date{}
\maketitle
%\noindent\textit{Email:} anamitroappu@gmail.com;\\\textit{Phone:} +91 9051798881;\\\textit{Website: }\url{https://sites.google.com/view/anamitro}

\section{Homotopy}
\noindent{}Is $\mathbb{R}\approx \mathbb{R}^2$? No, since removing one point from $\reals$ makes it disconnected, while removing a countable number of points from $\reals^2$ does not create a disconnection. But they are in many ways similar, and one is an expansion of the other. Homeomorphism is a very rigid condition imposed upon similarity, and does not reflect the necessary conditions for our working cases of similarity.

\begin{prob}We want a looser version of homeomorphism.
\end{prob}

\subsection{Homotopy of maps}
\begin{defn}
Let $f, g:X\tends Y$  be \textit{maps} (continuous functions). We say $f\simeq g$ if $\exists$ map $F:X\times I\tends Y~\ni~F(x,0)=f(x), F(x,1)=g(x)~\forall~x\in X$. $F$ is a homotopy from $f$ to $g$; $F:f\simeq g$.
\end{defn}

\begin{obs}``$\simeq$'' is an equivalence relation. Let $f\in\C(X, Y)$ be a map from $X$ to $Y$; $[f]=\lc g\in\C(X, Y)~\mid~f\simeq g\rc$, equivalence class of $f$.
\end{obs}

\textbf{Proof}. \begin{itemize}
  \item[(i)]For any map $f:X\tends Y$, $f\simeq f$.
  \item[] Here we simply consider $F(x,t)$ to be independent of $t$, and $F(x,t)=f(x)~\forall~x\in X$.
  \item[(ii)]Let $f, f\dash:X\tends Y$. Then, $f\simeq f\dash\Rightarrow f\dash\simeq f$.
  \item[]There exists homotopy $F:X\times I\rightarrow Y$ such that $F(x, 0)=f(x)$ and $F(x,1)=f\dash(x)$. We define $G:X\times I\rightarrow Y$ such that $G(x,t)=F(x, 1-t)~\forall~t\in[0,1]$. We have that $G$ is a homotopy between $f\dash$ and $f$.
  \item[(iii)]Let $f, f\dash, f\ddash:X\tends Y$. Then, $f\simeq f\dash,f\dash\simeq f\ddash\Rightarrow f\simeq f\ddash$.
  \item[]There exists $F:X\times I\tends Y$ such that $F(x,0)=f(x),~F(x,1)=f\dash(x)$ and $F\dash:X\times I\tends Y$ such that $F\dash(x,0)=f\dash(x),~F\dash(x,1)=f\ddash(x)$. Define $G:X\times I\tends Y$ with
$$G(x,t)=\begin{cases}
  F(x,2t)~&\text{for }t\in\left[0, \frac{1}{2}\right],\\
  F\dash(x,2t-1)~&\text{for }t\in\left[\frac{1}{2}, 1\right].
\end{cases}$$
We note that $G$ is well-defined and for $t=\dfrac{1}{2}$, $F(x,2t)=f\dash(x)=F\dash(x, 2t-1)$. Because $G$ is continuous on the two closed subsets $X\times\left[0, \dfrac{1}{2}\right]$ and $X\times\left[\dfrac{1}{2}, 1\right]$ of $X\times I$, by \emph{Pasting Lemma}, $G$ is continuous throughout and the required homotopy.\Square 
\end{itemize}


\subsection{Homotopy equivalence of spaces}
\begin{defn}
Suppose $\exists~f:X\tends Y,g:Y\tends X~\ni~g\circ f\simeq\id_X, f\circ g\simeq\id_Y$. Then we say that $X$ is homotopically equivalent to $Y$, $X\simeq Y$, or that $X$ and $Y$ have the same homotopy type.
\end{defn}

\begin{obs}
$X\approx Y\Rightarrow X\simeq Y$
\end{obs}

\textit{Proof.} $f^{-1}$ exists and is continuous. $g=f^{-1}$.

\begin{rem}
Converse is not true.
\end{rem}

\textit{Counterexample.} $X=\reals^2$, $Y=\lc (0,0)\rc$. $X\not\approx  Y$ since $X$ is not compact but $Y$ is. Let $F:\reals^2\times I\rightarrow\reals^2;~\lb x=\lb x_1, x_2\rb, t\rb\mapsto \lb 1-t\rb x~\left[x\in\reals^2, t\in[0,1]\right]$. Then, $F(x, 0)=\id_X$ and $F(x,1)=\id_Y$. So, there exists a homotopy from $\id_X$ to $\id_Y$.

Alternatively, we define $i:\lc (0,0)\rc\tends \reals^2, (0,0)\mapsto (0,0);~c:\reals^2\tends\lc (0,0)\rc.$ Then $i\circ c:\reals^2\tends\reals^2; x\mapsto (0,0)~\forall x\in\reals^2$ and $c\circ i:\lc (0,0)\rc\tends\lc (0,0)\rc; c\circ i=\id_Y$.\Square 

\begin{defn}
A path-connected space is called \emph{contractible} is it is homotopically equivalent to a one-point space, i.e., it has the same homotopy type as a point.
\end{defn}

\begin{obs}
Any 2 maps $f,g:X\tends Y$, where $Y$ is contractible, are homotopic.
\end{obs}

\textit{Proof.} Let Y be contractible. So there should exist a homotopy $H:1_Y\simeq f_c$ where $f_c:Y\tends \lc c\rc\subseteq Y$.

Let $X$ be another space and we consider a map $g:X\tends Y$. Then $H\circ g:X\times I\tends Y$ is a homotopy between $g$ and the constant map $\tilde{f}:X\tends \lc c\rc\subseteq Y$. So every map $g:X\tends Y$ is homotopic to the constant map $\tilde{f}$. Since homotopy is an equivalent relation, this implies that all maps $X\tends Y$ are homotopic.\Square

\begin{defn}
A \emph{retraction} is a map $r$ from a space $X$ to its subset $A$, such that $r|_A=\id_A$.

A deformation retraction is a homotopy $F:X\times I\tends A\subseteq X$ such that $F|_{A,I}=\id_A$. Here, additionally, the map $F$ is continuous. The retract is said to be ``relative $A$'' or $\rel A$.
\end{defn}

\subsection{Homotopy between Paths}
\noindent{}For a topological space $X$, a \emph{path} is a map $\s:[0,1]\tends X$. A path is called a \emph{loop} at $x\in X$ if $\s(0)=\s(1)=x$.

Let $\a$ and $\b$ be paths from $x$ to $y$ in $X$. We say that $\a$ \emph{is homotopic to} $\b$ relative to the endpoints $(0,1)$
$$\a\simeq\b\rel(0,1)$$
if $\exists~F:I\times I\tends X~\ni$
\begin{eqnarray*}
  F(s,0)=\a(s)~&\forall&s\in I;\\
  F(s,1)=\b(s)~&\forall&s\in I;\\
  F(0,t)=\a(0)=\b(0)=x~&\forall&t\in I;\\
  F(1,t)=\a(1)=\b(1)=y~&\forall&t\in I.
\end{eqnarray*}

We define $F_t:I\tends X;~F_t(s)=F(s,t)$ for each fixed $t\in I$. Therefore, $F_t$ is also a path from $x$ to $y$. Thus, homotopy can be seen as kind of a continuous deformation of path $\a$ to path $\b$.

We define $P\lb X;x,y\rb=\lc\a:I\tends X~|~\a(0)=x,a(1)=y\rc$. ``Being homotopic'' is an equivalence relation on $P\lb X;x,y\rb$ because

\begin{itemize}
  \item[(i)]Consider $F(s,t)=\a(s)\forall~s,t\in I$ which shows that $\a\simeq\a \rel(0,1)$, i.e., reflexivity.
  \item[(ii)]If $F:\a\simeq\b\rel(0,1)$ then $G:I\times I\tends X;~G(s,t)=F(s,1-t)$ gives $G:\b\simeq\a\rel(0,1)$. Symmetry.
  \item[(iii)]If we have $F:\a\simeq\b\rel(0,1)$ and $G:\b\simeq\g\rel(0,1)$, then we define the \emph{concatination} of $F$ and $G$ by
  $$(F*G)(s,t)=\begin{cases}
    F(s,2t)&\text{for }0\leq t\leq\frac{1}{2};\\
    G(s,2t-1)&\text{for }\frac{1}{2}\leq t\leq 1.
  \end{cases}$$
  which is continuous by the pasting lemma, whence $F*G:\a\simeq\g\rel(0,1)$. Transitivity.
\end{itemize}

\begin{note}
Concatination refers to basically traversing the first path in half the time with double speed, and then the second path in the other half of the time with double speed.
\end{note}

\subsection{Fundamental Group}
\noindent{}The equivalence class of $\a\in P(X;x,y)$ is denoted by $\ls\a\rs$ by $\ls\a\rs=\lc\a\dash\in P(X;x,y)~|~\a\simeq\a\dash\rel(0,1)\rc$. The set of all such equivalence classes is denoted by $P\ls X; x,y\rs$. Let's divert our discussion to loops at a point, say, $x_0$. We denote $P\ls X; x_0,x_0\rs$ by $\pi_1\lb X,x_0\rb\coloneq \lc\ls\s\rs~|~\s\text{ is a loop at }x_0\in X\rc$. We define a binary operation $[*]:\pi_1\lb X,x_0\rb\times\pi_1\lb X,x_0\rb\tends\pi_1\lb X,x_0\rb;~[\s][*][\t]\mapsto[\s *\t]$, where concatination $\s *\t:I\tends X$ is defined in our old way:
$$\s *\t(s)=\begin{cases}
  \s(2s)&\text{when }0\leq s\leq\frac{1}{2};\\
  \t(2s-1)&\text{when }\frac{1}{2}\leq s\leq 1.
\end{cases}$$
The operation $[*]$ is well-defined, i.e., it does not depend on the choice of representatives $\s$ and $\t$. If $\s\simeq\s\dash\rel(0,1)$ and $\t\simeq\t\dash\rel(0,1)$ then $\s*\t\simeq\s\dash *\t\dash\rel(0,1)$. Now we observe that
\begin{itemize}
\item[(i)]Of course, concatination gives another loop, and hence its equivalence class gives an element of $\pi_1\lb X, x_0\rb$.
\item[(ii)]$*$ is associative and hence so is $[*]$.
\item[(ii)]$[e]$ is the unit in $\pi_1\lb X, x_0\rb$ where $e:I\tends X;~e(s)=x_0~\forall~s\in I$, the \emph{constant loop} at $x_0$.
\item[(iii)]We observe that $\s\inv(s)=\s(1-s); [\s]\inv=\ls\s\inv\rs$ where $\s * \s\inv\simeq \s\inv * \s\simeq e\rel(0,1)$.
\end{itemize}
We conclude that $\pi_1\lb X, x_0\rb$ forms a group with $[*]$ as a binary operation, known as the \emph{fundamental group}.

If $x_1\in X$ is another point in $X$ then $\pi_1\lb X, x_1\rb$ makes sense.

\begin{thm}
If $X$ is path-connected then a path $\a$ from $x_0$ to $x_1$ induces an isomorphism $\a_*$ from $\pi_1\lb X, x_0\rb$ to $\pi_1\lb X, x_1\rb$ by
\begin{eqnarray*}
  \a_*:\pi_1\lb X, x_0\rb&\tends&\pi_1\lb X, x_1\rb;\\
   \ls\s\rs&\mapsto&\ls\a\inv *\s *\a\rs;\text{ describe $\a\inv$, loop $\s$, return to $x_1$}.
\end{eqnarray*}
\end{thm}
\noindent{}To check that $\a_*$ is a well-defined isomorphism, we need to show that
\begin{itemize}
  \item[(i)]$\a_*\lb \ls \s\rs\ls\t\rs\rb=\a_*\lb \ls \s\rs\rb\a_*\lb \ls\t\rs\rb;$
  \item[(ii)]$\lb\a_*\rb\inv=\lb\a\inv\rb_*$.
\end{itemize}
Because the operation $*$ is well-defined, if $\s$ is a loop based at $x_0$, then $\a\inv *\lb\s *\a\rb$ is a loop based at $x_1$. To show $\a_*$ to be a group isomorphism, we need
\begin{eqnarray*}
  \a_*\lb \ls \s\rs\rb\a_*\lb \ls\t\rs\rb&=&\lb\ls\a\inv * \s * \a\rs\rb\lb\ls\a\inv * \t * \s\rs\rb\\
  &=&\ls\a\inv * \s * \a *\a\inv * \t * \s\rs\\
  &=&\ls\a\inv\rs\ls\s *\t\rs\ls\a\rs\\
  &=&\a_*\lb \ls \s\rs\ls\t\rs\rb.
\end{eqnarray*}
Let $\b$ denote the path $\a\inv$, then $\b_*$ is the inverse for $\a_*$. For each $\ls\s\rs\in\p_1\lb X, x_1\rb$, $\b_*\lb\ls \s\rs\rb=\ls\b\inv *\s *\b\rs=\ls\a *\s *\a\inv\rs\Rightarrow\a_*\lb\b_*\ls\s\rs\rb=\ls\a\inv * \a *\s *\a\inv *\a\rs=\ls \s\rs$. Similarly $\b_*\lb\a_*\ls \s\rs\rb=\ls \s\rs$ for all $\ls \s\rs\in\p_1\lb X, x_0\rb$.

Moreover, if $\a\simeq\b\rel (0,1)$ for $\a,\b\in P\ls X; x_0, x_1\rs$, then $\a_*=\b_*$.

Thus, it makes sense to talk about the fundamental group of a space, irrespective of the base point.

\subsection{Homotopy groups}
Let $f:X\tends Y$, and $[f]$ be the equivalence class containing $f$, called its \emph{homotopy class}. The collection of all homotopy classes $X\tends Y$ is denoted by $[X,Y]$. In the set $\p_n(X)=\ls S^n, X\rs$, we can similarly define a binary operation for $f,g:S^n\tends X$,

$$\lb fg\rb\lb s_1,\dots, s_n\rb=\begin{cases}
  f\lb 2s_1, s_2,\dots, s_n\rb,~s_1\in\ls 0, \frac{1}{2}\rs,\\
  g\lb 2s_1-1, s_2,\dots, s_n\rb,~s_1\in\ls \frac{1}{2}, 1\rs.
\end{cases}$$
The identity will be the class of the constant map to the base point, and $-f\lb s_1, s_2,\dots, s_n\rb=f\lb 1-s_1, s_2,\dots, s_n\rb$. This group $\p_n(X)$ is called the \emph{$n$th (ordinary) homotopy group} of $X$. In fact, for $n\geq 2$ they are all abelian.


\subsection{Homotopy extension property}

\section{Cell complexes}
\subsection{Weak Topology}
\noindent{}Let $X$ be a topological space and let $W_\a, \a\in A$ be subsets of $X$ that in themselves are topological spaces. A weak topology on $X$ is the largest topology such that the inclusion map $W_\a\hookrightarrow X$ is continuous for each $\a\in A$. Equivalently, a subset $U$ of $X$ is open iff $U\cap W_\a\subseteq W_\a$ is open for each $\a\in A$. The collection $\lc U\subseteq X~|~ U\cap W_\a\subseteq_{\text{open}} W_\a~\forall~\a\in A\rc$ forms a topology on $X$, and every other topology that has the inclusion maps continuous, is contained in this weak topology.

Given $\lc W_\a\rc_{\a\in A}$ we can, if they are not disjoint, consider homeomorphic spaces that are disjoint, and take disjoint union $\displaystyle\displaystyle\coprod_{\a\in A} W_\a=X$, on which we can consider the weak topology.

\subsection{Identification maps}

\subsection{Adjunction spaces}\noindent{}Let $X$ and $Y$ be unbased topological spaces. We consider $W\subset X$ and let $\f:W\tends Y$ be a free map, i.e., a map which need not map base point to base point. On $Y\sqcup X$ we define the equivalence relation $w\sim \f(w)$. This gives us the identification space $Y\cup_\f X$, called the \emph{adjunction space}.

\subsection{Wedge}\noindent{}Let $\lb X_\a, x_\a\rb; \a\in A$ be based spaces. When we quotient the disjoint union $\displaystyle\coprod X_\a$ by $\lc x_\a~|~\a\in A\rc$, we basically identify all the base points into one, and this gives us the \emph{wedge} of $\lc \lb X_\a, x_\a\rb\rc_{\a\in A}$, denoted as $\displaystyle\bigvee_{\a\in A}X_\a$.

\subsection{CW complexes}
Note that for each $\a\in A$, there is an injection $i_\a:X_\a\tends \displaystyle\bigvee_{\a\in A}X_\a; x\mapsto \ls x\rs~\forall~x\in X_\a$. Furthermore, the maps $f_\a:X_\a\tends Y$ determine a unique map $f:\displaystyle\bigvee_{\a\in A}X_\a\tends Y~\ni~fi_\a=f_\a$.

If we consider $X$ to be an unbased space and assume $\f_\a:S^{n-1}=S_\a^{n-1}\tends X$ to be free maps for $\a\in A$, these determine a free map (by pasting lemma) $\f:\displaystyle\coprod_{\a\in A}S_\a^{n-1}\tends X$. Let $E_\a^{n-1}=E^{n}$ for all $\a\in A$. Since $\displaystyle\coprod_{\a\in A}S_\a^{n-1}\subset \displaystyle\coprod_{\a\in A}E_\a^{n}$, we can form the adjunction space by quotienting with respect to $\displaystyle\coprod_{\a\in A}S_\a^{n-1}$ as: $X\cup_{\a\in A}\displaystyle\coprod_{\a\in A}E_\a^{n}$. This procedure is called \emph{attaching $n$-cells}.

We consider an unbased Hausdorff topological space $X$, which we construct as follows: we consider a chain of subsets $X^0\subseteq X^1\subseteq X^2\subseteq \dots\subseteq X^{n-1}\subseteq X^n\subseteq\dots\subseteq X$, which we construct by considering $X^0$ as a $0$-cell or consisting of discrete vertices, and attaching $n$-cells to $X^{n-1}$ to form $X^n$. We have to assume that for each $n$, there exist \emph{attaching maps} $\f_\b:S^{n-1}\tends X^{n-1};\b\in B$ and we thus form $X^n\coloneq X^{n-1}\cup_\f \displaystyle\coprod_{\b\in B}E_\a^{n}$. Note that with each such formation of the disjoint union $X^{n-1}\sqcup  \displaystyle\coprod_{\b\in B}E_\a^{n}$ and subsequent quotienting, for each $\b\in B$, we get a continuous \emph{characteristic function} $\F_\b:\lb E^n, S^{n-1}\rb\tends \lb X^n, X^{n-1}\rb$ and $\F\vert_{S_\b^{n-1}}=\f_\b$. The image $\F_\b\lb E_\b^{n}-S_\b^{n-1}\rb=e_\b^n$ is called an \emph{open $n$-cell} (although it need not be open in $X$, and $\F_\b\vert_{E_\b^{n}-S_\b^{n-1}}: E_\b^{n}-S_\b^{n-1}\tends e_\b^n$ is a homeomorphism. The closure of $e_\b^n$ is $\F_\b\lb E_\b^n\rb$ and the topological boundary $\partial e_\b^n=\F_\b\lb S_\b^{n-1}\rb$. Then $\F_\b:\lb E^n, S^{n-1}\rb\tends \lb \overline{e}_\b^n, \partial e_\b^n\rb$ is a continuous function of pairs. This closed set $\overline{e}_\b^n=\F\lb E^n\rb\subseteq X^n\subseteq X$ is called a \emph{(closed) $n$-cell} that is being attached. $X$ as a whole is endowed with the weak topology with respect to all the closed cells $\lc \overline{e}_\b^n\rc_{\b\in B,n\in\naturals\cup\lc 0\rc}$. Here, $X^n-X^{n-1}$ is a disjoint union of the $e_\b^n$ and we can write $\displaystyle X^n=X^{n-1}\cup\bigcup_{\b\in B}e_\b^n$ or $\displaystyle X^n=X^{n-1}\cup\bigcup_{\b\in B}\overline{e}_\b^n$.

A CW complex is called finite if it has finitely many cells, and is called finite-dimensional (dimension $N$) if there exists $N\in\naturals$ such that $X^{N-1}\neq X$ but $X^n=X$ for all $n\geq N$.

%\subsection{Examples}
% \begin{thebibliography}{99}
%\bibitem[B]{thesis}
 %\end{thebibliography}

\end{document}
